% SOURCE_AUTHORITY: sga5_fr_workpass.tex, lines 15235--15244 (LF-normalized unit).
% SOURCE_SHA256: 791F4EFFC5E02832D5D77ED03518C8156D6F07E4C8238B03545DB93D883FBB28
Para demonstrar os lemas 1 e 2, necessitaremos do resultado seguinte:

\begin{statement}{Lema 3}
\begin{enumerate}[label=(\roman*)]
\item Sejam $P$ e $Q$ dois $A$-módulos. Se $P$ é projetivo, para toda aplicação $A_0$-linear $f_0:T(P)\to T(Q)$ existe uma aplicação $A$-linear $f:P\to Q$ tal que $T(f)=f_0$.
\item Suponhamos que o ideal $I$ seja nilpotente. Para todo $A_0$-módulo projetivo $P_0$ existe um $A$-módulo projetivo $P$ tal que $T(P)\simeq P_0$. Se, além disso, $P'$ for um $A$-módulo plano tal que $T(P')\simeq P_0$, então se tem $P'\simeq P$.
\end{enumerate}
\end{statement}

A afirmação (i) é clara, pois, se $u_P$ (resp. $u_Q$) designa a aplicação canônica de $P$ sobre $T(P)$ (resp. de $Q$ sobre $T(Q)$), o composto $f_0\circ u_P$ é uma aplicação $A$-linear de $P$ em $T(Q)$ e $u_Q$ é sobrejetiva; como $P$ é projetivo, existe uma fatoração de $f_0\circ u_P$ através de $u_Q$: $f_0\circ u_P=u_Q\circ f$, onde $f$ é uma aplicação $A$-linear de $P$ em $Q$; então $T(f)=f_0$.
